% 组合数学（综述）
% license CCBYSA3
% type Wiki

本文根据 CC-BY-SA 协议转载翻译自维基百科\href{https://en.wikipedia.org/wiki/Combinatorics}{相关文章}。

组合数学是数学的一个分支，主要研究计数（counting）——既作为获得结果的方法，又作为研究有限结构某些性质的目的。它与数学的许多其他领域密切相关，并在从逻辑到统计物理、从进化生物学到计算机科学等众多领域中有广泛的应用。

组合数学以其所处理问题的广泛性而著称。组合问题出现在纯数学的许多分支中，尤其是在代数学、概率论、拓扑学与几何学中\(^\text{[1]}\) ，同时也出现在众多应用领域。  
历史上，许多组合问题往往是独立提出的，作为在某个数学背景下出现的特定问题的特设性解法（ad hoc solutions）。然而，在二十世纪后期，强大而一般的理论方法被发展出来，使得组合数学成为一门独立的数学分支\(^\text{[2]}\)。组合数学中最古老且最易于理解的部分之一是图论（graph theory），它本身就与其他数学领域有着大量自然的联系。在计算机科学中，组合数学被广泛用于推导算法分析中的公式与估计。
\subsection{定义}
组合数学的确切范围并没有被普遍认同\(^\text{[3]}\)。正如 H. J. Ryser 所指出的，由于组合数学跨越了众多数学分支，因此要给出一个精确的定义是困难的\(^\text{[4]}\)。如果我们从其所研究的问题类型来描述这一领域，那么组合数学主要涉及以下几个方面：
\begin{enumerate}
  \item 枚举（enumeration）：对特定结构的计数，这些结构在广义上被称为排列（arrangements）或配置（configurations），它们通常与有限系统相关；
  \item 存在性（existence）：研究满足特定条件的结构是否存在；
  \item 构造性（construction）：探讨如何以一种或多种方式构造出这些结构；
  \item 优化（optimization）：在多种可能的结构或解中寻找“最优”的一个，无论其意义为“最大”、“最小”或符合其他最优性标准。
\end{enumerate}
Leon Mirsky 曾指出：“组合数学是一系列彼此关联的研究，它们有共同点，但在目标、方法以及所达到的统一程度上却差异很大。”\(^\text{[5]}\)一种定义组合数学的方式，或许是通过描述它的各个分支（subdivisions）、它们所处理的问题以及所采用的技术。下文将采用这种方式来介绍组合数学。然而，从历史的角度看，是否将某些主题纳入组合数学的范畴，也存在一定的传统性和争议性\(^\text{[6]}\)。尽管组合数学主要关注有限系统，但其中的一些问题与方法也可以推广到无限（尤其是可数）离散系统的情形。
\subsection{历史}

基本的组合概念与枚举结果在古代世界的许多地方早已出现。最早记录在案的组合技巧可追溯到公元前16世纪的《莱因德纸草书》（Rhind Papyrus）中的第79题。这一问题涉及一种几何级数，与斐波那契（Fibonacci）关于“由若干个1和2的组合形成给定总和”的计数问题有相似之处\(^\text{[7]}\)。印度医师苏什鲁塔（Sushruta）在其著作《Sushruta Samhita》中指出，从六种不同味道中取一种、两种、……直到全部六种，共可形成 $2^{6}-1=63$ 种组合。希腊历史学家普鲁塔克（Plutarch）记载了斯多亚学派哲学家克律西波斯（Chrysippus，公元前3世纪）与天文学家喜帕恰斯（Hipparchus，公元前2世纪）之间的一场关于一个精妙枚举问题的争论，该问题后来被证明与施罗德–喜帕恰斯数（Schröder–Hipparchus numbers）有关\(^\text{[8][9][10]}\)。更早时，阿基米德（Archimedes，公元前3世纪）在其著作《Ostomachion》中可能已考虑过一个平面拼图问题中不同拼图配置数（configurations）的计算\(^\text{[11]}\)，而阿波罗尼乌斯（Apollonius）失传的部分著作中也可能包含类似的组合思想\(^\text{[12][13]}\)。

在中世纪时期，组合数学的研究主要在欧洲之外持续发展。印度数学家摩诃毗罗（Mahāvīra，约公元850年）给出了排列数与组合数的计算公式\(^\text{[14][15]}\)，这些公式可能早在公元6世纪时就已为印度数学家所知\(^\text{[16]}\)。犹太哲学家兼天文学家亚伯拉罕·伊本·以斯拉（Abraham ibn Ezra，约公元1140年）发现了二项式系数的对称性；而数学家利维·本·格尔松（Levi ben Gerson，亦称Gersonides）则在1321年给出了其闭合形式\(^\text{[17]}\)。展示二项式系数之间关系的算术三角形（arithmetical triangle）早在10世纪的数学论文中就已出现，它后来被称为帕斯卡三角形（Pascal’s triangle）。在中世纪的英格兰，钟学（campanology）的研究产生了当今所谓凯莱图（Cayley graphs）上的哈密顿环（Hamiltonian cycles）实例\(^\text{[18][19]}\)。

文艺复兴时期，伴随数学与科学的复兴，组合数学也焕发了新的生命力。帕斯卡（Pascal）、牛顿（Newton）、雅各布·伯努利（Jacob Bernoulli）与欧拉（Euler）的著作成为这一新兴领域的奠基之作。进入近代，J.J. Sylvester（19世纪末）与 Percy MacMahon（20世纪初）的工作奠定了枚举组合学（enumerative combinatorics）与代数组合学（algebraic combinatorics）的基础。与此同时，图论（graph theory）因与四色问题（four color problem）的联系而引发广泛关注。

在20世纪下半叶，组合数学经历了迅速的发展，促成了数十种新的学术期刊与会议的创立\(^\text{[20]}\)。这一增长部分得益于其与其他数学分支（如代数、概率论、泛函分析、数论等）之间的新联系与应用。这些跨界互动模糊了组合数学与数学及理论计算机科学其他领域之间的边界，同时也在某种程度上导致了该领域的局部分化。

